\chapter{Fazit}
\label{sec:fazit}

\section{Auswertung und Zusammenfassung}

Alle Punkte aus der Zielsetzungen konnten im Rahmen der Projektarbeit umgesetzt werden. Die Ergebnisse zeigen, dass es mit einfachen Mitteln wie einer USB-Webcam und sehr preiswerter Software möglich ist, bereits verwertbare Ergebnisse zu generieren. Ursprünglich war angedacht eine Open Source Lösung als Fundament für den Eye Tracker zu verwenden. Nach ausgiebiger Marktrecherche konnte ein passendes Softwarefundament ausfindig gemacht werden. Der Beam Eye Tracker ist zwar nicht Open Source aber die Daten die er durch seine \acrshort{api} bereitstellen kann, bieten genug Spielraum, um individuelle Lösungen damit bauen zu können. Eine Verschmelzung zwischen Carla und dem Eye Tracker war ohne größeren Hürden möglich. Die Daten aus dem Instance Segmentation Sensor können zusammen mit den Werten des Eye Trackers ausgewertet werden. Somit ist es möglich nicht nur die geschätzten Koordinaten des Eye Trackers auf der Leinwand auszugeben, sondern auch das betrachtete Objekt innerhalb der Carla Simulationsumgebung zu bestimmen.
Während der Implementierung des Eye Tracking Systems sind an manchen Stellen auch Probleme aufgetreten. Die Belichtung war vor allem während der Kalibrierung ein Problem. Durch verschiedene Ansätze und Experimente im Fahrsimulator konnten hier mit externer Belichtung gute Ergebnisse erzielt werden.
Ein anderes Problem waren die verfälschten Tracking Daten, verursacht durch die gekrümmte Leinwand. Durch umfassende Analyse des Problems, einige mathematische Ansätze und Experimente im Fahrsimulator, konnte das Problem stark reduziert werden. Die in Abbildung \ref{fig:imag12} und \ref{fig:imag13} dargestellten Ergebnisse zeigen, dass gute Ergebnisse erzielt werden konnten. 
Die Tracking Daten sind allerdings nicht perfekt und lassen sich durch äußere Einflüsse wie schlechte Belichtung, Bedienungsfehler bei der Kalibrierung, Veränderung der Sitzposition nach der Kalibrierung oder zwei Personen gleichzeitig im Fahrsimulator negativ beeinflussen. Der Eye Tracker könnte allerdings zum aktuellen Entwicklungsstand bereits eingesetzt werden. Anhand der gemessenen Performance kann behauptet werden, dass bereits zum aktuellen Stand nützliche Erkenntnisse während einer Sitzung im Fahrsimulator mit einer Testperson gewonnen werden können. Im Ausblick sollen weitere Schritte und mögliche Ideen präsentiert werden, mit welchen über diesen Low Budget Ansatz hinaus, noch bessere Ergebnisse erzeugt werden könnten.


\section{Ausblick}

Angeknüpft an diese Projektarbeit kann eine Verbesserung des Testaufbaues erfolgen. Dazu zählen eine Ausbesserung der Lichtverhältnisse im Fahrsimulator, eine fest verbaute Kamera und gegebenfalls eine Blende, sodass nur der Fahrer und nicht zusätzlich der Beifahrer von der Kamera erkannt wird. Ebenso kann die Eye Tracker Genauigkeit verbessert werden, indem die Korrektur für den gekrümmten Bildschirm durch einen potentiell noch besseren Optimierungsalgorithmus durchgeführt wird. Auch eine Optimierung für die y-Achse könnte ein weitere Ansatz sein. Bspw. könnte somit der statische Offset herausgerechnet werden, wenn die Person während einer Fahrt in vertikale Richtung die Sitzposition verändert. Eine Integration von künstlicher Intelligenz zur Bestimmung der genauen Blickposition könnte ebenfalls die Genauigkeit verbessern.

Abschließend lassen sich weiterführende Analysen einbauen, wie die zeitliche Betrachtung von Blickpositionen und einer Aufbereitung der Blickposition in eine dreidimensionale Umgebung, sodass die Fahrt und die gemessene Blickposition anschaulich ausgewertet werden können. 

\textcolor{red}{In den Anhang der Arbeit das Demo Video einfügen}

\clearpage

\printglossaries
\printbibliography
